\scanpage{11}
\uline{Théorème} (à prouver)

\textcircled{1} Toute carte cellulaire \inserted{$\ne\varnothing$} connexe et simplement connexe quasi régulière (i.e. dont tous les sommets comb. et toutes les faces ont même ordre) est régulière, et elle est complète (i.e. les sommets comb. d'ordre fini sont effectifs) sauf dans le cas d'une carte isomorphe à $(\mathbb R^2\supset\mathbb R)$ (dont la \unsure{complétion est|complétion|complétée est} la 2sphère munie d'un grand-cercle ponctué\ldots).

\textcircled{2} Une \unsure{Carte|carte} cellulaire \inserted{$\ne\varnothing$} conn. simpl. connexe régulière et complète (appelons \deleted{les} \inserted{une telle carte} \textquotedblleft cartes standart\textquotedblright) est déterminée, à isomorphisme près, par \inserted{le couple $(p,q)$} \inserted{où $p$ est} l'ordre des sommets comb., $q$ l'ordre des faces ($p,q\in\overline{\mathbb N^*}$, entiers $1\leq p,q\leq+\infty$). Le couple $(p,q)$ peut prendre toutes les valeurs \deleted{sauf} $\in\overline{\mathbb N^*}\times\overline{\mathbb N^*}$ \uline{sauf} $(1,\unsure{q|p})$ \inserted{et $(p,1)$ si $p\ne2$}. Appelons $(C_{p,q},r_{p,q})$ la carte standard épinglée (i.e. munie d'un biarc-repère $r_{p,q}$) correspondant à $(p,q)$, déterminée à \deleted{homéomorphisme} \inserted{isom. unique} près dans la catégorie isotopique des cartes ]

\textcircled{3}
\[
\operatorname{Aut}(C_{p,q})\xleftarrow{\sim}
\underset{\text{\unsure{groupe cartographique|gpe cartographique}}}{{\unsure{\boldsymbol{\Gamma}|\Gamma}}^{0}}/
(\rho_s^p,\rho_f^q)\overset{\text{déf}}{=}\boldsymbol{\Gamma}_{p,q}
\]
\deleted{\textcircled{\unsure{3|4}}}

\textcircled{4} \uline{Corollaire} Soit $\unsure{\mathscr C|C|\mathcal C}=(X\supset K\supset S)$ une carte cellulaire complète \inserted{connexe} quasi-régulière de type $(p,q)$. Alors le revêtement universel \inserted{$\unsure{\overline C|\widetilde C|\widetilde{\mathcal C}}$} est \deleted{\unsure{rég}} une carte standart, le groupe fondamen%
% archive sheet number: 10
